\subsection{Skill taxonomy}

Baihu covers 64 manipulation skills across 2989 tasks. Rather than treating tasks as isolated labels, the dataset organizes manipulation behaviors through a reusable set of skill-level action semantics. These skills describe both low-level motion primitives and higher-level manipulation interactions, including object acquisition, placement, handover, articulated object interaction, contact-based motion, tool and material manipulation, and navigation-related actions.

Table~\ref{tab:skill-taxonomy} summarizes the 64 skills according to seven action-semantic categories. This taxonomy provides an intermediate level between individual task names and raw action trajectories: a single task may involve multiple skills, and the same skill can appear across different objects, scenes, and robot embodiments. This structure helps characterize the manipulation diversity of Baihu beyond simple task count.

% Baihu v2.0 skill taxonomy table generated from materials/技能.csv.
% Required packages: booktabs, tabularx, adjustbox.

\begin{table}[t]
\centering
\caption{Skill taxonomy of Baihu v2.0. The 64 manipulation skills are grouped according to their action semantics during data collection and annotation.}
\label{tab:skill-taxonomy}
\footnotesize
\setlength{\tabcolsep}{4pt}
\renewcommand{\arraystretch}{1.15}
\begin{adjustbox}{max width=\textwidth}
\begin{tabularx}{\textwidth}{p{0.30\textwidth}p{0.08\textwidth}X}
\toprule
Category & Count & Skills \\
\midrule
Object acquisition and release & 9 & Grasp, Pick, Pick up, Pinch, Suction, Hold, Lift, Lower, Release \\
Transport, placement, and handover & 10 & Move, Place, Carry, Transfer, HandOver, TakeOver, Throw, Discard, Takeout, Remove \\
Spatial arrangement and organization & 8 & Alignment, Adjust, Straighten, Stack, Unstack, Pack, Ensure, Several Times \\
Articulated, container, and electrical interaction & 11 & Open, Close, OpenBox, CloseBox, DrawOut, Insert, Plug, Unplug, PressButton, Zip, Unzip \\
Contact-based motion primitives & 9 & Push, Pull, Press, Rotate, Flip, Stretch, Touch, Swipe, Shake \\
Tool, fluid, and material manipulation & 10 & Pour, Scoop, Wipe, Water, Brush, Stir, Stick, Tear, Fold, Unfold \\
Navigation, perception, and expressive actions & 7 & Navigate, Point, Scan, Check, Hang, Armretract, Wave \\
\midrule
\textbf{Total} & \textbf{64} & -- \\
\bottomrule
\end{tabularx}
\end{adjustbox}
\end{table}


\subsection{Robot platform composition}

Baihu contains data from 13 robot platforms, covering a diverse set of robot embodiments and data sources. Table~\ref{tab:platform-comp} summarizes the platform-level composition of the dataset in terms of task count, episode count, frame count, and total duration. This platform-level reporting provides transparency about the dataset source composition while keeping the main emphasis on overall scale, platform coverage, and standardized data organization.

\begin{table}[t]
\centering
\caption{Platform-level composition of Baihu. The table reports tasks, episodes, frames, and hours for each of the 13 robot platforms.}
\label{tab:platform-comp}
\scriptsize
\begin{tabular}{lrrrr}
\toprule
Robot platform & Tasks & Episodes & Frames & Hours \\
\midrule
ARX\_Loong & 3 & 842 & 2223107 & 20.58 \\
Astribot S1 & 100 & 29889 & 40932735 & 379.01 \\
Agilex Cobot Magic & 112 & 33098 & 55115692 & 510.33 \\
UR5e & 49 & 14805 & 22939047 & 212.40 \\
D-Wheel & 231 & 13485 & 36172626 & 334.93 \\
Franka & 9 & 2090 & 4037181 & 37.38 \\
Agibot GO-1 & 1512 & 140718 & 459379181 & 4253.51 \\
Fourier GR-2 & 319 & 96686 & 103615634 & 959.40 \\
Leju KUAVO & 123 & 29095 & 42586155 & 394.32 \\
QingLong & 167 & 46506 & 91909479 & 851.01 \\
Tianji & 6 & 1575 & 2727875 & 25.26 \\
Galaxea R1 & 239 & 70111 & 96863491 & 896.88 \\
Agibot A2 & 119 & 34675 & 69847611 & 646.74 \\
\textbf{Total} & \textbf{2989} & \textbf{513575} & \textbf{1028349814} & \textbf{9521.75} \\
\bottomrule
\end{tabular}
\end{table}

The platform composition in Table~\ref{tab:platform-comp} shows that Baihu integrates robot operation data from multiple platforms with different embodiments, control interfaces, and task settings. This coverage supports unified data standardization and cross-platform model evaluation while preserving platform-specific state and action schemas.

\subsection{Scenario and task coverage}

Beyond scale and embodiment diversity, Baihu is designed to cover a broad range of manipulation scenarios and task structures. The dataset is intended to support robot learning in both daily-life and production-oriented environments. Representative scenario families include industrial manufacturing, household service, catering or food handling, retail or pharmacy-like environments, and other generalization scenarios. These scenario categories provide a practical organization for analyzing how robot policies transfer across environments and object distributions.

At the task level, Baihu contains 2989 tasks. These tasks include both short-horizon primitive behaviors and longer-horizon task chains. Many tasks can be viewed as compositions of lower-level manipulation skills, such as grasping, placing, pushing, pulling, handing over, inserting, pressing, opening, closing, sorting, and arranging. This composition makes the dataset useful not only for direct action prediction, but also for studying how low-level manipulation skills combine into longer task sequences.

\subsection{Object and interaction diversity}

Baihu contains robot interactions with diverse real-world objects. The objects span household items, kitchen utensils, retail goods, packages, industrial parts, tools, flexible materials, and irregularly shaped objects. These objects differ in size, shape, material, rigidity, surface texture, and manipulation affordance. Such object diversity is important for learning policies that do not overfit to a small set of familiar objects.

The dataset also includes different interaction patterns between robots and objects. Some trajectories involve single-object pick-and-place behaviors, while others involve object sorting, tool use, assembly-like manipulation, container interaction, drawer or door operation, and multi-object organization. This diversity supports the study of generalization across objects, tasks, and physical interaction types.

\subsection{Temporal horizon and skill composition}

Robot manipulation tasks in Baihu span multiple temporal horizons. Short-horizon trajectories typically correspond to local operations such as grasping, pressing, scanning, or simple placement. Medium-horizon trajectories may include object transfer, opening and closing, sorting, or simple multi-step rearrangement. Long-horizon trajectories involve more extended task chains, such as cleaning, assembly, loading or unloading, and multi-stage object organization.

This temporal diversity is important for embodied foundation model training. Short trajectories provide dense examples of low-level visuomotor control, while longer trajectories expose models to task phases, intermediate goals, action rhythm, and sequential dependencies. As a result, Baihu can support both low-level action prediction and higher-level behavior modeling.

\subsection{Data format}

Baihu is converted from HDF5 source records into the LeRobot v2.1 format. Each episode is represented as a trajectory consisting of temporally ordered observation-action pairs. Depending on the original robot platform and task setup, each trajectory may include visual observations, robot state information, action commands, task labels, and metadata. The standardized format is intended to support direct loading by robot learning pipelines and large-scale model training systems.

A typical Baihu trajectory can be represented as a task instruction or task label, synchronized visual observations, robot proprioceptive states, robot actions, episode metadata, and frame-level timestamps or ordering information.
